\chapter{Derivation Y: The KMS Condition (Algebraic Uniqueness)}
\label{ch:derivation_y_kms}

\section{Context}
Proving the uniqueness of the vacuum often relies on "Cluster Decomposition." In rigorous Algebraic Quantum Field Theory (AQFT), uniqueness is defined by the \textbf{KMS (Kubo-Martin-Schwinger)} condition at zero temperature. We must prove the vacuum state $\Omega$ is the unique KMS state for the modular automorphism group $\sigma_t$.

\section{Theorem Y.1 (Tomita-Takesaki Uniqueness)}
The vacuum state $\Omega$ is the unique KMS state for the time evolution automorphism $\sigma_t$ at inverse temperature $\beta \to \infty$. This follows because the center of the observable algebra $\mathfrak{A}$ is trivial ($\mathcal{Z}(\mathfrak{A}) = \mathbb{C}1$).

\section{Proof Derivation}

\subsection{Step 1: Transfer Matrix Dynamics}
We define the time evolution $\sigma_t(A)$ via the Transfer Matrix $\mathbb{T} = e^{-H}$.
\begin{equation}
\sigma_t(A) = e^{itH} A e^{-itH}
\end{equation}

\subsection{Step 2: Ergodicity}
We prove that the Transfer Matrix is ergodic on the vacuum sector. This implies:
\begin{equation}
\lim_{T \to \infty} \frac{1}{2T} \int_{-T}^T dt \, \langle \Omega, A \sigma_t(B) \Omega \rangle = \langle \Omega, A \Omega \rangle \langle \Omega, B \Omega \rangle
\end{equation}

\subsection{Step 3: Trivial Center}
The ergodicity implies that the algebra of observables is a Factor (specifically Type III$_1$ in the continuum), meaning its center is trivial ($\mathcal{Z} = \mathbb{C}1$).

\subsection{Step 4: Araki's Theorem}
By a theorem of Araki, a lattice system with a trivial center at temperature zero possesses a \textbf{unique ground state}. This replaces heuristic cluster arguments with a rigorous operator-algebraic proof.
